\section{Week 1: Introduction to Interaction Design}

\subsection{Topic Summary}

\begin{keybox}[Learning Objectives]
By the end of this week, you will be able to:
\begin{itemize}
  \item \textbf{HAI1}: Identify key human-AI interaction concepts and discuss how these relate to human agency and initiative within the resulting system.
  \item \textbf{HAI2}: Describe the human-AI interaction design process and its distinction from conventional human-centred design.
\end{itemize}
\end{keybox}

\begin{keybox}[Big Picture]
This week introduces the foundations of interaction design: what it means to design for people using technology, and the theoretical concepts of \term{affordances} and \term{signifiers} that underpin good design. It also introduces \term{user-centred design} processes and how generative AI is changing the ``design rules'' we knew.

The key insight: \textbf{you are already a designer}---you naturally think about how objects should work. This course teaches you to apply that intuition systematically.
\end{keybox}

\begin{keybox}[What Examiners Like to Ask]
\begin{itemize}
  \item Define interaction design and its goals; apply to a given screenshot/system.
  \item Distinguish affordances from signifiers with examples.
  \item Identify false/hidden affordances in a design and suggest fixes.
  \item Explain the Double Diamond model and when you diverge/converge.
  \item Describe how AI differs as a ``design material'' from traditional technologies.
  \item Given a context, propose interview questions to understand user needs.
\end{itemize}
\end{keybox}

\subsection{Overview + Core Concepts + Extended Theory}

\subsubsection*{What is Interaction Design?}

\begin{defbox}[Interaction Design]
``The design of interactive products to support the way people communicate and interact in their everyday and working lives.'' (Textbook definition)
\end{defbox}

\paragraph*{Scope of Interaction Design}
Interaction design is \textbf{anything involving technology and people}:
\begin{itemize}
  \item Websites, mobile apps, voice assistants, touchless interfaces, AR, full-body interaction
  \item Also: data visualization, architecture, industrial design, car dashboards, border control flows
  \item All subsumed by \term{user experience} (UX)---designing interactions is part of the broader UX
\end{itemize}

\paragraph*{Three Eras of Interaction Design}
\begin{center}
\renewcommand{\arraystretch}{1.2}
\begin{tabular}{@{} l p{9cm} @{}}
\toprule
\textbf{Era} & \textbf{Characteristics} \\
\midrule
Traditional ID & Deterministic systems; clear rules; 10--15 year history \\
Deep Learning AI & Designer selects data, trains models; bias/fairness/trust critical \\
Generative AI & Foundation models outside your control; fine-tuning only; rules still emerging \\
\bottomrule
\end{tabular}
\end{center}

\paragraph*{Interaction Design Goals}
\begin{enumerate}
  \item \textbf{Easy to learn}: Users can start using it without extensive training.
  \item \textbf{Effective to use}: Accomplishes the intended task efficiently.
  \item \textbf{Enjoyable user experience}: Or at least not dreadful!
\end{enumerate}

\begin{exbox}[Good vs Bad Design: Parking Signs]
\textbf{Bad}: A sign listing parking restrictions in complex text (``No parking Mon--Fri 7am--9am except...''). Hard to parse while driving.

\textbf{Good}: Visual representation showing time blocks when parking is allowed/forbidden. Easy to scan at a glance.

Both communicate the same information, but the good design is \textit{easy to learn}, \textit{effective}, and less frustrating.
\end{exbox}

\paragraph*{Five Characteristics of Design}
\begin{enumerate}
  \item \textbf{Changing situations}: Design involves shaping/deploying artifacts to change a situation.
  \item \textbf{Exploring possible futures}: Imagining better alternatives to current problems.
  \item \textbf{Framing problem + solution in parallel}: Understanding what's wrong while proposing fixes.
  \item \textbf{Thinking through sketching}: Using tangible representations (prototypes, sketches).
  \item \textbf{Addressing multiple aspects}: Instrumental, technical, aesthetical, and \textbf{ethical} throughout.
\end{enumerate}

\begin{tipbox}[The Uncomfortables]
Visit \texttt{theuncomfortable.com}---objects that look functional but fail in use (wellies with sandal straps, fork with unusable handle). These illustrate how our brains \textbf{automatically simulate interactions} and detect design failures.
\end{tipbox}

\subsubsection*{Affordances}

\begin{defbox}[Affordance]
``The perceived and actual properties of the thing, primarily those fundamental properties that determine just how the thing could possibly be used.'' (Norman, 1988)

Affordances describe the \textbf{relationship between object and user}---what actions are possible.
\end{defbox}

\paragraph*{Key Insight}
Affordances are about \textbf{functionality enabled by the object}:
\begin{itemize}
  \item Button $\to$ affords pushing
  \item Lever $\to$ affords moving up/down
  \item Knob $\to$ affords rotating
  \item Handle $\to$ affords pulling/gripping
\end{itemize}

\paragraph*{Affordance Matrix}

\begin{center}
\renewcommand{\arraystretch}{1.3}
\begin{tabular}{@{} l | c c @{}}
\toprule
& \textbf{Affordance Present} & \textbf{Affordance Absent} \\
\midrule
\textbf{Perceptible} & \textcolor{ThmGreen}{Good} (can do, can see) & \textcolor{ThmGreen}{Good} (can't do, can see) \\
\textbf{Not Perceptible} & \textcolor{TipRed}{Hidden affordance} & \textcolor{ThmGreen}{Good} (can't do, don't see) \\
\bottomrule
\end{tabular}
\end{center}

\begin{itemize}
  \item \textbf{Perceptible affordance}: Good---user knows what's possible.
  \item \textbf{Hidden affordance}: Bad---action is possible but user doesn't know.
  \item \textbf{False affordance}: Bad---user thinks action is possible but it isn't.
\end{itemize}

\begin{exbox}[The Door Problem]
A door with a handle on both sides:
\begin{itemize}
  \item From pull side: handle $\to$ ``pull'' (correct affordance).
  \item From push side: same handle $\to$ user tries to pull, but door only pushes.
\end{itemize}
\textbf{Solution}: Push side should have flat plate (no handle), pull side has handle. The affordance matches the action.
\end{exbox}

\paragraph*{Affordances in Digital Systems}
Physical affordances are often visible (you see a handle). Digital affordances are \textbf{hidden}---you have pixels, but which are clickable?
\begin{itemize}
  \item Text box $\to$ affords typing
  \item Button $\to$ affords clicking
  \item Dropdown $\to$ affords selecting
\end{itemize}
Digital systems rely heavily on \term{signifiers} to make affordances visible.

\subsubsection*{Signifiers}

\begin{defbox}[Signifier]
``Signifiers specify how people discover those possibilities: signifiers are signs, perceptible signals of what can be done.'' (Norman)

The perceivable part of an affordance is a signifier. If deliberately placed by a designer, it's a \term{social signifier}.
\end{defbox}

\paragraph*{Affordance vs Signifier}
\begin{center}
\renewcommand{\arraystretch}{1.3}
\begin{tabular}{@{} l p{5cm} p{5cm} @{}}
\toprule
& \textbf{Affordance} & \textbf{Signifier} \\
\midrule
Definition & What actions are possible & How users discover those actions \\
Example (button) & Can tap (touch target exists) & Visual button shape, label, icon \\
Example (door) & Can push/pull & Handle (pull) or plate (push) \\
\bottomrule
\end{tabular}
\end{center}

\begin{exbox}[ChatGPT Interface]
\begin{itemize}
  \item \textbf{Affordance}: Can type text (text box), can send (button), can select suggestions (clickable items).
  \item \textbf{Signifiers}: Text box outline, arrow icon on send button, hover highlighting on suggestions.
\end{itemize}
\end{exbox}

\begin{tipbox}[Norman's Design Principle]
``When simple things need pictures, labels, or instructions, the design has failed.''

Complex things may need explanation, but simple things (doors!) should be self-evident through good affordances and signifiers.
\end{tipbox}

\paragraph*{The Remote Control Example}
A grandmother's remote was ``simplified'' by covering unused buttons with tape, leaving only the essential ones visible. This is a \textbf{signifier intervention}---reducing cognitive load by hiding irrelevant affordances.

\subsubsection*{User-Centred Design (UCD)}

\begin{defbox}[User-Centred Design]
``An approach to systems design and development that aims to make interactive systems more usable by focusing on the use of the system and applying human factors/ergonomics and usability knowledge and techniques.'' (ISO 9241-210:2019)

Key: Design teams involve users throughout the process via research and design techniques.
\end{defbox}

\paragraph*{The Double Diamond Model}

\begin{center}
\textbf{Diamond 1: Doing the Right Thing} $\to$ \textbf{Diamond 2: Doing Things Right}
\end{center}

\begin{enumerate}
  \item \textbf{Discover} (diverge): Explore the problem space, understand context/users.
  \item \textbf{Define} (converge): Synthesize insights into a \term{design brief}.
  \item \textbf{Develop} (diverge): Ideate solutions, prototype multiple options.
  \item \textbf{Deliver} (converge): Refine, test, implement the final solution.
\end{enumerate}

\begin{tipbox}[Key Insight]
You may spend 60--70\% of time just figuring out \textbf{what} to build (Diamond 1). The design brief---understanding the real problem---is itself a critical contribution.
\end{tipbox}

\paragraph*{Iteration is Expected}
\begin{itemize}
  \item Not waterfall: you rarely have all information upfront.
  \item Arrows go back: prototyping may reveal you misunderstood the problem.
  \item Similar to Agile: iterate, test assumptions, learn by doing.
\end{itemize}

\paragraph*{Stanford d.school / IDEO 5 Stages}
\begin{enumerate}
  \item \textbf{Empathize}: Understand users and their needs.
  \item \textbf{Define}: Frame the problem clearly.
  \item \textbf{Ideate}: Generate many possible solutions.
  \item \textbf{Prototype}: Build quick, testable versions.
  \item \textbf{Test}: Get feedback, iterate.
\end{enumerate}

\paragraph*{The ``Faster Horse'' Tension}
\begin{quote}
``If I had asked people what they wanted, they would have said a faster horse.'' (attributed to Henry Ford)
\end{quote}
\textbf{Users know their needs but not the solutions.} Ford noticed people wanted ``a faster horse that doesn't get tired''---understanding the \textit{need} (speed, endurance) led to the car \textit{solution}.

This is the core UCD tension: extract needs from users, but don't expect them to design the solution.

\subsubsection*{AI as Design Material}

\begin{keybox}[AI Changes the Rules]
Generative AI brings the first new UI paradigm in 60 years. Traditional HCI rules may not fully apply.
\end{keybox}

\paragraph*{Material Friction}
From the blog ``AI as Design Material'' (Haiyan Zhang, Microsoft):
\begin{quote}
``Each medium has a life of its own and resists what you want to do. Clay resists, wood splinters... and AI \textbf{misinterprets}.''
\end{quote}

\textbf{Good design happens in understanding and working with material friction}---the quirks, resistance, and unexpected behaviors of your medium.

\paragraph*{The Design Squiggle for AI}
\begin{itemize}
  \item Traditional design: Messy at start $\to$ converges to clarity $\to$ execution.
  \item AI design: The squiggle may \textbf{never fully converge}---systems change underneath you.
\end{itemize}

\begin{exbox}[Case Study: Micro-Narratives]
\textbf{Problem}: Researchers needed young people to share difficult social media experiences at scale, but open-text surveys yield poor responses.

\textbf{Design Brief}: Create a remote narrative elicitation method that is:
\begin{itemize}
  \item Open-ended (not a survey)
  \item Consistent across participants
  \item Low cognitive burden
\end{itemize}

\textbf{Solution}: Three sequential LLM calls:
\begin{enumerate}
  \item \textbf{Scaffold questions}: Ask structured questions (What happened? Context? Feelings? Worst part?).
  \item \textbf{Extract answers}: Pull key phrases from responses.
  \item \textbf{Craft story}: Combine extractions into coherent vignette.
\end{enumerate}

\textbf{Result}: In a study (N=254), micro-narratives were:
\begin{itemize}
  \item 4.5$\times$ more likely to be ``easier to capture experience''
  \item 6.0$\times$ more likely to be ``helpful for making sense of experience''
  \item Statistically significantly less difficult, more accurate
\end{itemize}

\textbf{Key design choices}:
\begin{itemize}
  \item Show 3 story variations (different emotional tones) so user doesn't just accept first output.
  \item Let user adapt/edit the final story.
  \item Template-based: change questions $\to$ adapt to any domain.
\end{itemize}
\end{exbox}

\subsubsection*{Worked Examples}

\begin{exbox}[Identify Affordances and Signifiers]
\textbf{Scenario}: A coffee mug with a sliding lock mechanism.

\textbf{Analysis}:
\begin{itemize}
  \item \textbf{Affordances}: Button (can press), sliding lock (can slide).
  \item \textbf{Signifiers}: Icons showing locked/unlocked states, button sized for thumb.
\end{itemize}
\end{exbox}

\begin{exbox}[Redesign a Bad Affordance]
\textbf{Scenario}: Glass door with identical handles on both sides (push/pull unclear).

\textbf{Problem}: False affordance---handle suggests pull, but one side requires push.

\textbf{Redesign}:
\begin{itemize}
  \item Push side: Flat plate (affords pushing, not gripping).
  \item Pull side: Handle (affords gripping and pulling).
  \item No labels needed if affordances are correct.
\end{itemize}
\end{exbox}

\begin{exbox}[Apply Double Diamond]
\textbf{Context}: You're asked to design an app for NHS vaccination passports.

\textbf{Diamond 1 (What to build)}:
\begin{itemize}
  \item Discover: Interview event organizers, NHS staff, travelers. What are their pain points?
  \item Define: Design brief---e.g., ``A mobile app that lets event staff verify vaccination status in <10 seconds without internet.''
\end{itemize}

\textbf{Diamond 2 (How to build it)}:
\begin{itemize}
  \item Develop: Sketch multiple solutions (QR code, NFC, biometric).
  \item Deliver: Prototype top 2, test with users, iterate to final design.
\end{itemize}
\end{exbox}

\subsubsection*{Key Definitions Summary}

\begin{center}
\renewcommand{\arraystretch}{1.3}
\begin{tabular}{@{} l p{9cm} @{}}
\toprule
\textbf{Term} & \textbf{Definition} \\
\midrule
Interaction Design & Designing interactive products to support communication/interaction in daily life \\
Affordance & Perceived/actual properties determining how something could be used \\
Signifier & Perceptible signal of what actions are possible \\
User-Centred Design & Iterative design focusing on users and their needs throughout \\
Double Diamond & Diverge/converge twice: discover$\to$define, develop$\to$deliver \\
Design Brief & Document defining what problem to solve and constraints \\
\bottomrule
\end{tabular}
\end{center}

\subsubsection*{Recommended Reading}
\begin{itemize}
  \item \textbf{Textbook}: Sections 1.3--1.5 (Introduction to Interaction Design)
  \item \textbf{Blog}: ``AI as Design Material'' (Haiyan Zhang)---treats AI as a medium with resistance
  \item \textbf{Article}: ``Signifiers, Not Affordances'' (Don Norman)---clarifies the distinction
  \item \textbf{TED Talk}: Tim Brown on design thinking
\end{itemize}


